\documentclass[12pt]{article}

\input(paper/preamble)

% ---------------------------------------------
% ---------------- TITLE ----------------------
% ---------------------------------------------
\title{Crowded Primaries and Weakened Nominees: Behavioral Distortion in a Two-Stage Electoral Contest}
\author{Katherine Parslow}
\date{Last Updated: \today}


\begin{document}
\maketitle

\section*{Model \& Notation}

\subsection*{Setup}

\subsection*{Primary Elections}

\subsection*{General Election}

\subsection*{Payoffs}

\subsection*{Timing}

\begin{itemize}
    \item $N \ge 2$ candidates compete for party nomination.
    \item Candidates have valence $v_i \in \R$ drawn i.i.d. from distribution $F$ with strictly positive density $f$ on $\R$. The vector $\vprof = \{v_1,...,v_N\}$ is drawn at the start of the game and is observed by all candidates.
    \item Candidates choose whether to take action $a_i = \{0,1\}$, which boosts their relative standing in the primary election but harms their viability in the next-stage general election.
    \item Candidates are ranked according to their primary index
    \[
        q_i^P = v_i + a_i + \varepsilon_i
    \]
   where $\varepsilon\sim$ EV1 are i.i.d.\ candidate-level shocks. The nominee, denoted $w$, is determined by the candidate with the highest primary index. Integrating over $\varepsilon_i$ results in a probability of nomination for candidate $i$ given by
   \[
    \Pr(w = i \mid \vprof, \aprof) = \frac{\ex{v_i + a_i}}{\sum_{j=1}^N \ex{v_j + a_j}}
   \]
   \item The nominee $w$ advances to the general election where their general-election index is
   \[
        q_i^{GE} = v_i - \delta\ a_i + \eta_i
   \]
   where $\eta_i \sim N(0,1)$ is an idiosyncratic shock. Note that primary and general election shocks are independent of each other. I normalize the out-party opponent's index to $0$ and the probability of nominee $w$ winning the general election is
   \[
        \Pr(\GE\ win \mid v_w, a_w) = \Pge{v_w}{\delta}{a_w}
   \]
   \item Candidates are office-motivated and receive payoffs of $1$ from winning office and $0$ otherwise. Because holding office requires winning both the primary and general elections, candidate $i$'s expected utility is given by
   \begin{align*}
        u_i(\vprof, \aprof) &= \Pr(w = i \mid \vprof, \aprof)\cdot \Pr(\GE\ win \mid v_i, a_i)\\
        &= \frac{\ex{v_i + a_i}}{\sum_{j=1}^N \ex{v_j+a_j}} \cdot \Pge{v_i}{\delta}{a_i} 
   \end{align*}
   \item The timing of the game is as follows:
   \begin{enumerate}
       \item Nature draws the vector of valences $\vprof = (v_1,...,v_N)$;
       \item Candidates observe $\vprof$ and choose $a_i \in \{0,1\}$ for all $i$;
       \item The primary election outcome is realized and the party nominee is selected;
       \item The general-election outcome is realized and payoffs are determined.
   \end{enumerate}
\end{itemize}

\end{document}